\Subhead{3.4.1}{Функция Грина и интеграл Дюамеля. Шум в колебательном контуре. Фильтрация шума}

Причинная линейная стационарная система с нулевыми начальными условиями имеет отклик
\begin{equation}
 y(t)=\int_{-\infty}^th(t-s)x(s)\dd s,\quad h(t<0)=0,
 \qquad H(\omega)=\int h(t)\e^{\ii\omega t}\dd t.
\end{equation}
$h$ — реакция на дельта-импульс, или функция Грина с учётом принятой нормировки входа. Если $g(t)$ — реакция на единичную ступень, то $h=g'$. Формула Дюамеля может быть записана через $g$ и производную входа: $y(t)=g(t)x(0)+\int_0^tg(t-s)\dot x(s)\dd s$ при соответствующем включении в $t=0$.

Ковариация выхода содержит двойную свёртку; преобразование Фурье превращает её в
\begin{equation}
 S_y(\omega)=|H(\omega)|^2S_x(\omega),\qquad
 \sigma_y^2=\frac1{2\pi}\int|H|^2S_x\dd\omega.
\end{equation}
Устойчивость фильтра и существование интеграла необходимы: идеальный интегратор на белом шуме даёт нестационарное блуждание, а не стационарный выход с конечной дисперсией.

Для $m\ddot x+b\dot x+kx=F(t)$
\begin{equation}
 H_{xF}=\frac1{k-m\omega^2-\ii b\omega},\quad
 S_x=\frac{S_F}{(k-m\omega^2)^2+b^2\omega^2}.
\end{equation}
Контур выделяет узкую полосу около собственной частоты: белый вход становится узкополосным выходом. При тепловом $S_F=2bk_BT$ интегрирование даёт $\langle x^2\rangle=k_BT/k$, $\langle\dot x^2\rangle=k_BT/m$. Уменьшение потерь повышает вершину спектра и сужает его, сохраняя равновесную площадь. Для RC-фильтра $H=(1-\ii\omega RC)^{-1}$ эффективная односторонняя полоса белого шума $\int_0^\infty|H(2\pi f)|^2\dd f=1/(4RC)$, а не частота среза $1/(2\pi RC)$.
